\section{The classification problem}

A finite nonempty set $A$ in a Euclidean space is \emph{Ramsey} if, for
every integer $r\ge2$, there is an integer $D\ge1$ such that every map
$c:\R^D\to\{1,\ldots,r\}$ is constant on a congruent copy of $A$.
Congruence means preservation of every pairwise distance, at the original
scale. The coloring is an arbitrary function. In particular, no
measurability or other regularity is assumed.

Erd\H{o}s, Graham, Montgomery, Rothschild, Spencer and Straus introduced
Euclidean Ramsey theory and proved that every finite Ramsey set is
spherical~\cite[Theorem~13]{egm-rss}. Graham later conjectured that
sphericity is also sufficient~\cite{graham-survey}. Major positive
results established the Ramsey property for nondegenerate triangles
and then all nondegenerate simplices, by Frankl and
R\"odl~\cite{frankl-rodl-triangles,frankl-rodl-simplices}; finite sets
with a soluble transitive isometry group and cyclic trapezoids, by
K\v{r}\'i\v{z}~\cite{kriz-permutation,kriz-trapezoids}; and vertex sets
of regular polytopes, by Cantwell~\cite{cantwell-polytopes}.

Leader, Russell and Walters proposed a different characterization:
a finite set is Ramsey exactly when it embeds in a finite transitive
Euclidean set~\cite[Conjecture~A]{lrw-transitive}. Transitivity means
that the isometries of the finite set act transitively on its points.
An isometric subset of such a set is called \emph{subtransitive}.
Their group-theoretic and Hales--Jewett formulations led to the
Block Sets Conjecture; subsequent work studies the possible block
sizes~\cite{ilw-blocks}. The link with soluble symmetry has also
expanded: Karamanlis embedded every simplex in a finite product of
regular polygons~\cite{karamanlis-tori}, and Behague constructed
soluble transitive extensions for many known Ramsey classes and
further permutation configurations~\cite{behague-subsoluble}.
Ivan, Leader and Walters developed generalized-prism constructions,
with Ramsey conclusions under a soluble symmetry
hypothesis~\cite{ilw-prisms}. Moore and, independently, Mirabi proved
that adjoining a point outside the affine span of a Ramsey set
preserves the Ramsey property~\cite{moore-pyramid,mirabi-extensions}.
The affine-span restriction leaves the case of arbitrary coplanar
circle points beyond this extension theorem.

P\'alv\"olgyi has recently given a seven-point circle configuration
that is not Ramsey, disproving the spherical
conjecture~\cite[Theorem~1]{palvolgyi-heptagon}.
The appendix to that version also states that generic seven-point circle
configurations are not Ramsey. Its author's note announces a proof that
all finite transitive sets are Ramsey and says that its exposition is
still being prepared.\footnote{P\'alv\"olgyi attributes the writing, proofs, and
ideas in that manuscript to ChatGPT. He labels the additional appendix
claims as unchecked by him.}
Our main theorem gives a necessary and sufficient algebraic condition
for an arbitrary finite Euclidean configuration.

\subsection{The exact tensor criterion}

The problem is to characterize these finite sets up to congruence.
A singleton is Ramsey. For any other set, let $d\ge1$ be its affine
dimension, and choose a congruent representative
$A=\{a_1,\ldots,a_s\}\subset\R^d$ whose affine span is $\R^d$.
The points are distinct and $s\ge2$. Define
\[
 p_i=\begin{pmatrix}1\\a_i\end{pmatrix},\qquad
 F=\Q\bigl((a_i)_\alpha:1\le i\le s,\ 1\le\alpha\le d\bigr),
 \qquad B=F\otimes_{\Q}F.
\]
The commutative ring $B$ has a multiplication homomorphism
\[
 m_F:B\longrightarrow F,\qquad m_F(x\otimes y)=xy.
\]
An equality in $B$ is stronger than the equality obtained by applying
$m_F$; this map need not be injective. Index the entries of $p_i$ and
the rows and columns of a $(d+1)$-by-$(d+1)$ matrix by $0,1,\ldots,d$.
Thus index $0$ is the constant coordinate.

\begin{theorem}[Classification]\label{thm:classification}
Let $A=\{a_1,\ldots,a_s\}\subset\R^d$ be a finite set of at least two
points with affine span $\R^d$, and define $p_i,F,B,m_F$ as above.
Then $A$ is Ramsey if and only if there is a matrix
$P\in\operatorname{Mat}_{d+1}(B)$ such that
\begin{align}
 (p_i\otimes1)^T P(1\otimes p_i)&=0
       &&(1\le i\le s),\label{eq:field-evaluations}\\
 m_F(P_{\alpha\beta})&=\delta_{\alpha\beta}
       &&(1\le\alpha,\beta\le d).\label{eq:field-spatial}
\end{align}
Here the entries of $p_i\otimes1$ and $1\otimes p_i$ are formed
componentwise. There are no separate conditions on the multiplied
constant row or constant column of $P$.
\end{theorem}

The condition is independent of the chosen congruent representative
in its affine dimension; we also prove this directly through its tensor
formulation. It uses exact relations in the coordinate field. It is
not a procedure for recovering those relations from numerical
approximations to arbitrary real coordinates.

\subsection{Geometric consequences}

Applying $m_F$ to the evaluation equations gives an ordinary sphere
equation
\[
 \|a_i\|^2+\ell\cdot a_i+c=0\qquad(1\le i\le s)
\]
for some $\ell\in F^d$ and $c\in F$. The tensor condition requires the
evaluation identities to hold before multiplication. This is the
distinction between sphericity and the Ramsey property that the
classification detects.

Every subtransitive set satisfies the tensor condition. Another
sufficient condition is that evaluation at the points of a spherical
set gives independent functionals on the space of quadratic
polynomials. It applies to every nonempty set of at most five points
on a circle, and hence to every cyclic quadrilateral. In
Corollary~\ref{cons:kites}, we combine this result with the
non-subtransitive cyclic kites of Leader, Russell and
Walters~\cite[Corollary~2]{lrw-kites}. Those kites are Ramsey, disproving
the necessity direction of their proposed characterization and their
